\section{Discussion}\label{sec:discussion}

The outcomes indicate that one correlation-graph PF-STGT can produce coupled day-ahead load and stress forecasts from mutual abstractions on Bangladesh national data. Following we examine the key patterns and what they represent for this grid.

\subsection{Correlation Graphs Versus Geography}

Geographical adjacency is intuitive- divisions that share a border may switch power-but it handles every border link equally regardless of how significantly loads in practice co-move. On Bangladesh national footprint, evening-peak demand correlations are positive across divisions, governed by shared growth, seasonality, and holidays rather than border proximity. At $\tau{=}0.65$, the correlation graph maintains 33 of 36 pairs; geographical adjacency maintains only 21 border links and overlooks divisions that move jointly without exchanging a border.

A5 (geographical graph) slightly strengthens OSI MAE (0.0340) but increases demand error by 9.33~MW over S2, which indicates that the optimal graph for stress is not inherently the best for load. Geographical conditioning considerably degrades demand ($p{=}1.48{\times}10^{-4}$) because it imposes signal passing through a sparse, equal-weight skeleton. Correlation weighting rather provides graded edges that match predictive coupling and can be enhanced furthermore by graph awareness (Eq.~\eqref{eq:spatial}). Similar behaviour has been reported on other national-scale systems where synchronisation dominates local exchange \cite{Zhao2024GraphAttentionTransformer,Bentsen2023GraphTransformer}.

\subsection{Multi-Task Learning Trade-offs}

Demand and OSI draw on the same calendar, limitation, and generation context but answer different questions: regional megawatt levels versus a single bounded stress score. Shared encoding lets the model pick up limitation-related structure---also visible in attribution for G8 (limitation stack) and G6 (calendar/trend)---that a demand-only model would have to recover indirectly. Without the loss rescaling in Eq.~\eqref{eq:loss}, demand gradients dominate and the OSI head barely trains.

Demand-only A4 reaches 86.89~MW, 1.76~MW below S2, because the stress head adds a mild regularising pull away from pure demand fitting. Joint training still provides useful OSI output ($R^2{=}0.745$). In implementation, a centre that needs only load forecasts may prefer A4; one that requires paired load and stress insights can receive that small demand cost for much superior stress predictions than B07 or conventional baselines supplied.

\subsection{Bangladesh Grid Context}

Three features of the dataset form these observations. Nine divisions reserve, report shedding, and limitation in the same daily records, which enables composite stress supervision feasible without additional telemetry. Dhaka dominates national demand ($\sim$35.7\%) and error, so macro metrics downplay how responsive planning is to the capital region; that concentration emerges throughout all model families and appears more similar to a data-scale challenge than a model failure. Lastly, daily reporting with a seven-day lookback matches administrative practice, which assists describe why shallow recurrent and T-GCN models lag full transformer capability on a extensively linked correlation graph.

\subsection{Deployment Notes}

A single operation could revise regional evening-peak trajectories and national OSI scalar from one seven-day input window every day. Operators might examine both outputs during day ahead planning, watch limitation stack, and report Dhaka separately and calendar inputs flagged by attribution. We would expect shadow-mode deployment alongside current tools first, with division-level MAE validates grounded in offline accuracy, not proof of enhanced grid consistency or economic consequence in the field.

\subsection{Limitations and Future Work}

All observations come from one Bangladesh timeline at daily precision with internal standards only; there is no external replication or live implementation trial. Statistical testing includes demand MAE; OSI gains lack hypothesis test or intervals. Attributes indicate model behavior, not causation, and demand attributions differ with permutation rankings ($\rho{=}{-}0.564$). Uncertainty quantification and whole S2 residual diagnostics were not finished.

Natural next steps comprise repeating the protocol on similar South Asian systems, adding probabilistic heads, testing adaptive correlation graphs, and running operator in the loop studies of dual task explanation monitoring interfaces.
